\item \textbf{ACTIVIDAD: Simulación estadística del decaimiento radiactivo.}
Esta actividad simula la desintegración estadística de núcleos radiactivos utilizando dados comunes.
\begin{enumerate}
    \item Primero, plantee el siguiente procedimiento mentalmente: coloque $100$ dados en una bolsa y agite durante unos segundos. Lance los dados en una mesa. Cada lanzamiento representará un intervalo de tiempo $\Delta t$. Retire todos los dados que muestren el número $1$ en su cara superior y regístrelos a un lado. Registre el número restante $N$ de dados. Vuelva a poner los dados restantes en la bolsa y repita sucesivamente hasta que queden pocos dados. Pronostique teóricamente la vida media del procedimiento: el número de lanzamientos después de los cuales queda exactamente la mitad de los dados.
    \item Finalmente, realice la actividad físicamente. Registre los resultados, grafique el logaritmo natural del número $N$ de dados que quedan después de cada lanzamiento contra el número de lanzamientos $n$, determine experimentalmente la vida media y compárela con su predicción teórica.
\end{enumerate}